% 专题七：内积与柯西不等式
\section{专题七：内积与柯西不等式}

\begin{learnobjective}
\begin{itemize}
  \item 会用平方非负证明二维柯西不等式，并说明等号条件；
  \item 会用内积、长度与投影解释柯西不等式；
  \item 会用判别式或配方法给出一般形式的证明；
  \item 会通过构造向量解决平方和、分式和与条件最值问题。
\end{itemize}
\end{learnobjective}

\begin{preview}
\subsection{导语：同一个问题，五个看法}
$(a^2+b^2)(c^2+d^2)$ 与 $(ac+bd)^2$ 谁大？
平方非负能证（$(ad-bc)^2\geq0$），向量的长度与夹角能证，二次函数判别式能证，配方也能证。
同一个柯西不等式，把代数、几何、函数三个领域连了起来。
本专题用五个视角把三个“种子问题”各做一遍。

\subsection{种子问题}
先自己做，再看后面的五个视角。
\begin{enumerate}
  \item 证明二维柯西不等式：$(a^2+b^2)(c^2+d^2)\geq(ac+bd)^2$，并说明等号条件。
  \item 证明：对 $a,b,c>0$，$(a+b+c)\left(\dfrac1a+\dfrac1b+\dfrac1c\right)\geq9$。
  \item 已知 $x+y=1$，求 $F(x,y)=\dfrac{x^2}{2}+\dfrac{y^2}{3}$ 的最小值。
\end{enumerate}
\end{preview}

\begin{explore}
\subsection{视角一：来源——从平方非负到柯西}

\key{要点}
\begin{itemize}
  \item 出发点仍是平方非负：$(ad-bc)^2\geq0$。
  \item 展开即得二维柯西不等式。
  \item 等号条件：$ad-bc=0$，即两向量 $(a,b)$ 与 $(c,d)$ 线性相关。
\end{itemize}

\begin{example}
证明：$(a^2+b^2)(c^2+d^2)\geq(ac+bd)^2$。
\end{example}
\begin{solution}
作差：
\[
  (a^2+b^2)(c^2+d^2)-(ac+bd)^2
  =a^2d^2+b^2c^2-2abcd
  =(ad-bc)^2\geq0.
\]
等号当且仅当 $ad=bc$（即两向量 $(a,b)$、$(c,d)$ 线性相关）。
\key{（回头看：种子问题 1）}
\end{solution}

\begin{problem}
证明：$(a^2+b^2+c^2)(x^2+y^2+z^2)\geq(ax+by+cz)^2$。
\end{problem}
\begin{solution}
$(a^2+b^2+c^2)(x^2+y^2+z^2)-(ax+by+cz)^2=(ay-bx)^2+(az-cx)^2+(bz-cy)^2\geq0$。
\end{solution}

\subsection{视角二：几何——内积、长度与投影}

\key{要点}
\begin{itemize}
  \item 向量内积：$\langle\mathbf u,\mathbf v\rangle=\sum a_ib_i$；长度：$\|\mathbf u\|=\sqrt{\sum a_i^2}$。
  \item 几何意义：$\langle\mathbf u,\mathbf v\rangle=\|\mathbf u\|\,\|\mathbf v\|\cos\theta$，而 $|\cos\theta|\leq1$，故
  \[
    |\langle\mathbf u,\mathbf v\rangle|\leq\|\mathbf u\|\,\|\mathbf v\|.
  \]
  \item 等号当且仅当两向量线性相关（方向相同或相反，或其一为零向量）。
\end{itemize}

\begin{example}
用内积证明 $(a^2+b^2)(c^2+d^2)\geq(ac+bd)^2$。
\end{example}
\begin{solution}
取 $\mathbf u=(a,b)$，$\mathbf v=(c,d)$，则
$\langle\mathbf u,\mathbf v\rangle=ac+bd$，$\|\mathbf u\|^2=a^2+b^2$，$\|\mathbf v\|^2=c^2+d^2$。
由内积不等式即得。
\key{（回头看：种子问题 1 的几何证法）}
\end{solution}

\begin{problem}
用内积解释：为什么 $(a^2+b^2)(c^2+d^2)=(ac+bd)^2$ 当且仅当 $(a,b)$ 与 $(c,d)$ 共线？
\end{problem}
\begin{solution}
等号即 $|\cos\theta|=1$，也就是两向量方向相同或相反（夹角 $0$ 或 $\pi$），即线性相关。
\end{solution}

\subsection{视角三：函数——判别式与配方}

\key{要点}
\begin{itemize}
  \item 关键表达式：$\|\mathbf u-t\mathbf v\|^2\geq0$ 对任意实数 $t$ 成立。
  \item 展开得关于 $t$ 的二次函数：
  \[
    \|\mathbf v\|^2t^2-2\langle\mathbf u,\mathbf v\rangle\,t+\|\mathbf u\|^2\geq0.
  \]
  \item 它恒非负 $\iff$ 判别式 $\Delta\leq0$，即柯西不等式。
  \item 也可以配方：最小值在 $t=\dfrac{\langle\mathbf u,\mathbf v\rangle}{\|\mathbf v\|^2}$ 处取得。
\end{itemize}

\begin{example}
用判别式证明 $n$ 维柯西不等式：
$\left(\sum a_ib_i\right)^2\leq\left(\sum a_i^2\right)\left(\sum b_i^2\right)$。
\end{example}
\begin{solution}
设 $f(t)=\left(\sum a_i^2\right)t^2-2\left(\sum a_ib_i\right)t+\left(\sum b_i^2\right)$，
则 $f(t)=\sum(a_it-b_i)^2\geq0$ 对任意 $t$ 成立。
故
\[
  \Delta=4\left(\sum a_ib_i\right)^2-4\left(\sum a_i^2\right)\left(\sum b_i^2\right)\leq0,
\]
即得。
\key{（回头看：种子问题 1 的函数证法）}
\end{solution}

\begin{problem}
把 $f(t)=\sum(a_it-b_i)^2$ 配方，求其最小值，并说明它与柯西不等式的关系。
\end{problem}
\begin{solution}
\[
  f(t)=\left(\sum a_i^2\right)\left(t-\frac{\sum a_ib_i}{\sum a_i^2}\right)^2
  +\frac{\left(\sum a_i^2\right)\left(\sum b_i^2\right)-\left(\sum a_ib_i\right)^2}{\sum a_i^2}.
\]
最小值非负，故 $\left(\sum a_ib_i\right)^2\leq\left(\sum a_i^2\right)\left(\sum b_i^2\right)$。
\end{solution}

\subsection{视角四：应用——构造向量}

\key{要点}
\begin{itemize}
  \item 柯西不等式的常用形式（Engel 形式）：对 $a_i>0$，
  \[
    \sum_{i=1}^n\frac{x_i^2}{a_i}\geq\frac{(x_1+\cdots+x_n)^2}{a_1+\cdots+a_n}.
  \]
  \item 关键不是计算，而是\key{构造向量}：把目标式写成内积与长度的关系。
  \item 先看目标式的结构（平方和？分式和？乘积和？），再决定如何拆。
\end{itemize}

\begin{example}
证明：对 $a,b,c>0$，$(a+b+c)\left(\dfrac1a+\dfrac1b+\dfrac1c\right)\geq9$。
\end{example}
\begin{solution}
取 $\mathbf u=(\sqrt a,\sqrt b,\sqrt c)$，$\mathbf v=\left(\dfrac1{\sqrt a},\dfrac1{\sqrt b},\dfrac1{\sqrt c}\right)$，
则 $\langle\mathbf u,\mathbf v\rangle=3$，$\|\mathbf u\|^2=a+b+c$，$\|\mathbf v\|^2=\dfrac1a+\dfrac1b+\dfrac1c$。
由柯西不等式 $3^2\leq(a+b+c)\left(\dfrac1a+\dfrac1b+\dfrac1c\right)$，即得。
等号当且仅当 $a=b=c$。
\key{（回头看：种子问题 2）}
\end{solution}

\begin{example}
已知 $x+y=1$，求 $F(x,y)=\dfrac{x^2}{2}+\dfrac{y^2}{3}$ 的最小值。
\end{example}
\begin{solution}
取 $\mathbf u=\left(\dfrac{x}{\sqrt2},\dfrac{y}{\sqrt3}\right)$，$\mathbf v=(\sqrt2,\sqrt3)$，
则 $\langle\mathbf u,\mathbf v\rangle=x+y=1$。
由柯西不等式：
\[
  1^2\leq\left(\frac{x^2}{2}+\frac{y^2}{3}\right)(2+3),
\]
故 $F(x,y)\geq\dfrac15$。
等号要求 $\dfrac{x}{2}=\dfrac{y}{3}$，结合 $x+y=1$ 得 $x=\dfrac25$，$y=\dfrac35$。
\key{（回头看：种子问题 3）}
\end{solution}

\begin{problem}
已知 $x+y=1$（$x,y>0$），求 $\dfrac1x+\dfrac1y$ 的最小值。
\end{problem}
\begin{solution}
$(x+y)\left(\dfrac1x+\dfrac1y\right)\geq(1+1)^2=4$，故 $\dfrac1x+\dfrac1y\geq4$，等号当 $x=y=\dfrac12$。
\end{solution}

\begin{problem}
已知 $x+y+z=1$，求 $\dfrac1x+\dfrac4y+\dfrac9z$ 的最小值。
\end{problem}
\begin{solution}
$(x+y+z)\left(\dfrac1x+\dfrac4y+\dfrac9z\right)\geq(1+2+3)^2=36$，
故最小值为 $36$，等号当 $x=\dfrac16$，$y=\dfrac13$，$z=\dfrac12$。
\end{solution}

\begin{problem}
已知 $a^2+b^2=1$，求 $a+b$ 的最大值。
\end{problem}
\begin{solution}
由柯西不等式，$(a+b)^2\leq(1^2+1^2)(a^2+b^2)=2$，故 $a+b\leq\sqrt2$。
等号当 $a=b=\dfrac{\sqrt2}{2}$ 时成立。
\end{solution}
\key{（贯穿题：$a^2+b^2$ 型已在专题二、六出现）}

\subsection{视角五：运算——与凸性的联系}

\key{要点}
\begin{itemize}
  \item 平方范数是凸函数：对 $0\leq t\leq1$，
  \[
    \|t\mathbf x+(1-t)\mathbf y\|^2\leq t\|\mathbf x\|^2+(1-t)\|\mathbf y\|^2.
  \]
  \item 加权恒等式（$w_i>0$）：
  \[
    \left(\sum w_i\right)\left(\sum w_ix_i^2\right)-\left(\sum w_ix_i\right)^2
    =\sum_{i<j}w_iw_j(x_i-x_j)^2\geq0.
  \]
  它同时解释了加权柯西、加权 Jensen 与“方差非负”。
  \item 二次型：$Q(x,y)=ax^2+2bxy+cy^2$。当 $a>0$ 时配方得
  $Q(x,y)=a\left(x+\dfrac ba y\right)^2+\left(c-\dfrac{b^2}{a}\right)y^2$，
  故 $Q\geq0$ 对所有 $x,y$ 成立 $\iff ac-b^2\geq0$。
\end{itemize}

\begin{example}
用加权恒等式证明加权柯西不等式。
\end{example}
\begin{solution}
取 $w_i>0$，则
\[
  \left(\sum w_i\right)\left(\sum w_ix_i^2\right)-\left(\sum w_ix_i\right)^2
  =\sum_{i<j}w_iw_j(x_i-x_j)^2\geq0,
\]
即 $\left(\sum w_ix_i\right)^2\leq\left(\sum w_i\right)\left(\sum w_ix_i^2\right)$。
\end{solution}

\begin{problem}
用柯西不等式证明：对任意实数 $a,b$，$a^2+b^2\geq2ab$（提示：取 $\mathbf u=(a,b)$，$\mathbf v=(b,a)$）。
\end{problem}
\begin{solution}
$(a^2+b^2)(b^2+a^2)\geq(ab+ba)^2=4a^2b^2$，故 $(a^2+b^2)^2\geq4a^2b^2$，
即 $a^2+b^2\geq2|ab|\geq2ab$。
\end{solution}
\end{explore}

\begin{inclasstest}
\subsection{变式训练}
只改一个条件，重新做一遍。
\begin{enumerate}
  \item 证明：$(a^2+b^2+c^2)(x^2+y^2+z^2)\geq(ax+by+cz)^2$。
  \item 已知 $a,b,c>0$ 且 $a+b+c=1$，求 $\dfrac1a+\dfrac1b+\dfrac1c$ 的最小值。
  \item 已知 $x+2y=1$（$x,y>0$），求 $\dfrac1x+\dfrac1y$ 的最小值。
  \item 已知 $x+y+z=1$，求 $\dfrac1x+\dfrac4y+\dfrac9z$ 的最小值。
\end{enumerate}
\end{inclasstest}

\begin{afterclass}
\subsection{知识来源对照表}
\begin{center}
\footnotesize
\setlength{\tabcolsep}{4pt}
\renewcommand{\arraystretch}{1.35}
\begin{tabular}{p{3.4cm}p{4.0cm}p{5.6cm}}
\toprule
结论 / 方法 & 最早从哪来 & 后面从哪个角度看 \\
\midrule
平方非负 & 视角一（来源） & 视角三（函数）、视角五（运算） \\
内积与投影 & 视角二（几何） & 视角三（函数） \\
判别式 / 配方 & 视角三（函数） & 视角二（几何） \\
构造向量（Engel 形式） & 视角四（应用） & 视角二（几何） \\
加权恒等式 & 视角五（运算） & 视角四（应用） \\
二次型非负 & 视角五（运算） & 视角一（来源） \\
\bottomrule
\end{tabular}
\end{center}

\key{下一步}：三个核心支柱——函数与凸性（专题六）、内积与二次型（本专题）、集合与逻辑（专题一）——
可以互相解释：平方范数既是凸函数，又是内积诱导的长度。
\end{afterclass}
